\documentclass[10pt]{beamer}

\usepackage[ngerman]{babel}
\usepackage{amsmath, amssymb}

\usetheme{Boadilla}

\author[\url{https://fdf-uni.github.io/ft}]{}
\title{Tutorium 2}
\subtitle{\texorpdfstring{Funktionentheorie\vspace*{-1.5cm}}{Funktionentheorie}}
\date{05. Mai 2026}

\newcommand{\iu}{\mathrm{i}}
\renewcommand{\Re}{\operatorname{Re}}
\renewcommand{\Im}{\operatorname{Im}}

\begin{document}
\begin{frame}
	\titlepage
\end{frame}
\begin{frame}
	\frametitle{Kurven}
	\pause
	Für Definition einer Kurve, ihrer Umkehrung und Eigenschaften wie \emph{geschlossen}, \emph{einfach} oder \emph{äquivalent}, s. Vorlesung (im Wesentlichen ist eine Kurve eine stetige Abbildung $z \colon [a, b] \to \Omega \subset \mathbb{C}$ mit $z \in C_p^1([a, b])$).
\end{frame}
\begin{frame}
	\frametitle{Integration entlang von Kurven}
	\begin{definition}[Kurvenintegral]
		Gegeben einer Kurve $\gamma$ mit Parametrisierung $z \colon [a, b] \to \mathbb{C}$ und \glqq{}Knickstellen\grqq{} $a = a_0 < \ldots < a_K = b$ definieren wir das \emph{Integral von $f$ entlang $\gamma$} mittels
		\[
			\int_{\gamma} f(z) \,\mathrm{d}z := \sum_{k=1}^{K} \int_{a_{k-1}}^{a_k} f(z(t)) z'(t) \,\mathrm{d}t.
		\]
	\end{definition}
	\pause
	Merkregel für/Intuition zum Faktor $z'(t)$ (rein heuristisch):
	\pause
	\begin{itemize}
		\item \glqq{}Wie bei der Substitutionsregel\grqq{}
		      \pause
		\item Bezieht \glqq{}Geschwindigkeit\grqq{} der Parametrisierung mit ein
		      \pause
		\item Unabhängigkeit von Parametrisierung (vorausgesetzt, parametrisierte Kurven sind äquivalent im Sinne der Definition aus der Vorlesung)
	\end{itemize}
\end{frame}
\begin{frame}
	\frametitle{(Holomorphe) Stammfunktionen}
	\begin{definition}
		Sei $\Omega \subset \mathbb{C}$ offen und $f \colon \Omega \to \mathbb{C}$.
		Dann heißt $F \colon \Omega \to \mathbb{C}$ \emph{Stammfunktion (von $f$)}, falls $F$ holomorph in $\Omega$ ist mit $F' = f$.
	\end{definition}
	\pause
	\begin{theorem}
		Mit $\Omega$, $f$ und $F$ wie oben, gilt für jede Kurve $\gamma$ in $\Omega$ mit Anfangspunkt $w_0$ und Endpunkt $w_1$, dass
		\[
			\int_{\gamma} f(z) \,\mathrm{d}z = F(w_1) - F(w_0).
		\]
		Insbesondere gilt, für jede geschlossene Kurve $\gamma$ (also mit $w_0 = w_1$):
		\[
			\int_{\gamma} f(z) \,\mathrm{d}z = 0.
		\]
	\end{theorem}
\end{frame}
\end{document}
